\documentclass[14pt,a4paper,oneside]{extbook}
\usepackage[utf8]{vietnam}
\usepackage{amsmath,amssymb,mathrsfs,fancyhdr,enumerate,multirow,makecell}
\usepackage{tkz-euclide,tikz-3dplot,tikz,tkz-tab}
\usetikzlibrary{arrows.meta,calc,intersections,angles,quotes,backgrounds,shapes.geometric,patterns,through,shadows,positioning}
\usepackage{yhmath}
\usepackage[top=0.7cm, bottom=0.7cm, left=1.0cm, right=1.0cm, headheight=16pt, headsep=3mm, footskip=5mm]{geometry}
\usepackage[hidelinks,unicode]{hyperref}
\usepackage{qrcode}
\usepackage[many]{tcolorbox}
\usepackage{fontawesome}

\pagestyle{fancy}
\fancyhf{}
\fancyhead[L]{\footnotesize\bfseries\color{teal!85!black}\faBook~SỔ TAY KIẾN THỨC TOÁN 4}
\fancyhead[R]{\footnotesize\bfseries\color{blue!85!black}Lớp Toán Thầy Lĩnh -- loptoanthaylinh.vn}
\fancyfoot[C]{\footnotesize\bfseries\color{gray!80!black}-- Sổ tay kiến thức: Bài 64. Phép chia phân số --}
\renewcommand{\headrulewidth}{0.4pt}

\begin{document}

% --- BANNER TIÊU ĐỀ BÀI HỌC ---
\noindent
\begin{tcolorbox}[
  enhanced,
  colback=teal!8!white,
  colframe=teal!80!black,
  arc=3mm, boxrule=1pt,
  left=3mm, right=3mm, top=2mm, bottom=2mm,
  halign=center
]
  {\Large\bfseries\color{teal!90!black}\faBookmark\ BÀI 64: PHÉP CHIA PHÂN SỐ}\\[1.5pt]
  {\footnotesize\itshape\color{black!75} Phân số đảo ngược, quy tắc nhân đảo ngược và chia cho số tự nhiên trên 1 trang A4}
\end{tcolorbox}

\vspace{-2.5mm}

\noindent
\begin{minipage}[t]{0.49\textwidth}\vspace{0pt}
\begin{tcolorbox}[
  enhanced,
  colback=blue!4!white,
  colframe=blue!75!black,
  arc=2mm, boxrule=0.8pt,
  left=2mm, right=2mm, top=1.8mm, bottom=1.8mm,
  title={\footnotesize\bfseries\color{white}\faRefresh\ 1. Phân số đảo ngược}
]
{\footnotesize
\begin{itemize}\setlength{\itemsep}{1pt}\setlength{\leftskip}{-2.5mm}
  \item \textbf{Khái niệm:} Đảo ngược vị trí giữa tử số và mẫu số của một phân số ta được \textbf{phân số đảo ngược}.
  \item Phân số đảo ngược của $\dfrac{a}{b}$ là $\mathbf{\dfrac{b}{a}}$ ($a \ne 0, b \ne 0$).
  \item \textbf{Ví dụ:}
  \begin{itemize}\setlength{\itemsep}{0.5pt}\setlength{\leftskip}{-2.5mm}
    \item Đảo ngược của $\dfrac{3}{5}$ là $\dfrac{5}{3}$.
    \item Đảo ngược của $\dfrac{1}{4}$ là $\dfrac{4}{1} = 4$.
    \item Đảo ngược của số tự nhiên $5$ ($=\frac{5}{1}$) là $\dfrac{1}{5}$.
  \end{itemize}
\end{itemize}
}
\end{tcolorbox}
\end{minipage}
\hfill
\begin{minipage}[t]{0.49\textwidth}\vspace{0pt}
\begin{tcolorbox}[
  enhanced,
  colback=teal!4!white,
  colframe=teal!75!black,
  arc=2mm, boxrule=0.8pt,
  left=2mm, right=2mm, top=1.8mm, bottom=1.8mm,
  title={\footnotesize\bfseries\color{white}\faShareAlt\ 2. Quy tắc chia phân số}
]
{\footnotesize
Muốn chia một phân số cho một phân số, ta lấy \textbf{phân số thứ nhất nhân với phân số thứ hai đảo ngược}:
$$\dfrac{a}{b} : \dfrac{c}{d} = \dfrac{a}{b} \times \dfrac{d}{c} = \dfrac{a \times d}{b \times c}$$
\textbf{Chia cho số tự nhiên:}\\
$$\dfrac{a}{b} : c = \dfrac{a}{b \times c} \quad (c \ne 0)$$
\textbf{Số tự nhiên chia phân số:}\\
$$a : \dfrac{b}{c} = a \times \dfrac{c}{b} = \dfrac{a \times c}{b}$$
}
\end{tcolorbox}
\end{minipage}

\vspace{-2.5mm}
% --- KHỐI 3: SƠ ĐỒ TƯ DUY TIKZ ---
\noindent
\begin{tcolorbox}[
  enhanced,
  colback=orange!4!white,
  colframe=orange!80!black,
  arc=2mm, boxrule=0.8pt,
  left=2mm, right=2mm, top=1.5mm, bottom=1.5mm,
  title={\footnotesize\bfseries\color{white}\faSitemap\ SƠ ĐỒ TƯ DUY: TIẾN TRÌNH 3 BƯỚC PHÉP CHIA PHÂN SỐ}
]
\centering
\begin{tikzpicture}[
  box/.style={draw, rounded corners=2pt, align=center, font=\scriptsize\bfseries, inner sep=3pt, minimum height=1.1cm},
  arr/.style={-{Stealth[length=2.5mm]}, thick, shorten >=1.5pt, shorten <=1.5pt}
]
  \node[box, fill=purple!12, draw=purple!80!black, text width=3.4cm] (b1) at (0,0) {BƯỚC 1: GIỮ NGUYÊN\\Phân số bị chia $\mathbf{\dfrac{a}{b}}$\\Không thay đổi};

  \node[box, fill=blue!10, draw=blue!80!black, text width=3.4cm] (b2) at (4.6,0) {BƯỚC 2: CHUYỂN PHÉP TÍNH\\Đổi dấu $(:)$ thành $(\times)$\\Nhân với phân số đảo};

  \node[box, fill=teal!10, draw=teal!80!black, text width=3.6cm] (b3) at (9.4,0) {BƯỚC 3: ĐẢO NGƯỢC SỐ CHIA\\$\mathbf{\dfrac{c}{d} \longrightarrow \dfrac{d}{c}}$\\Tử $\leftrightarrow$ Mẫu hoán đổi};

  \draw[arr, purple!80!black] (b1) -- (b2);
  \draw[arr, blue!80!black] (b2) -- (b3);
\end{tikzpicture}
\end{tcolorbox}

\vspace{-2.5mm}
% --- KHỐI 4: VÍ DỤ & MẸO TRÁNH BẪY ---
\noindent
\begin{minipage}[t]{0.49\textwidth}\vspace{0pt}
\begin{tcolorbox}[
  enhanced,
  colback=green!4!white,
  colframe=green!60!black,
  arc=2mm, boxrule=0.8pt,
  left=2mm, right=2mm, top=1.8mm, bottom=1.8mm,
  title={\footnotesize\bfseries\color{white}\faCheckCircle\ Ví dụ thực hành chuẩn mực}
]
{\footnotesize
\textbf{1. Chia hai phân số:}\\
$\dfrac{3}{7} : \dfrac{2}{5} = \dfrac{3}{7} \times \dfrac{5}{2} = \dfrac{3 \times 5}{7 \times 2} = \mathbf{\dfrac{15}{14}}$.\\
$\dfrac{4}{9} : \dfrac{8}{3} = \dfrac{4}{9} \times \dfrac{3}{8} = \dfrac{4 \times 3}{9 \times 8} = \dfrac{1 \times 1}{3 \times 2} = \mathbf{\dfrac{1}{6}}$.\\
\textbf{2. Chia với số tự nhiên:}\\
$\dfrac{6}{7} : 3 = \dfrac{6}{7 \times 3} = \dfrac{2}{7}$; \quad $4 : \dfrac{2}{3} = 4 \times \dfrac{3}{2} = \mathbf{6}$.
}
\end{tcolorbox}
\end{minipage}
\hfill
\begin{minipage}[t]{0.49\textwidth}\vspace{0pt}
\begin{tcolorbox}[
  enhanced,
  colback=red!4!white,
  colframe=red!70!black,
  arc=2mm, boxrule=0.8pt,
  left=2mm, right=2mm, top=1.8mm, bottom=1.8mm,
  title={\footnotesize\bfseries\color{white}\faExclamationTriangle\ Mẹo tránh bẫy trắc nghiệm}
]
{\footnotesize
\begin{itemize}\setlength{\itemsep}{1.5pt}\setlength{\leftskip}{-2.5mm}
  \item \textbf{Bẫy đảo ngược nhầm phân số thứ nhất:}\\
  Chỉ đảo ngược \textbf{phân số thứ hai (số chia)}, tuyệt đối giữ nguyên phân số thứ nhất!
  \item \textbf{Bẫy lấy tử chia tử, mẫu chia mẫu:}\\
  Không được lấy tử chia tử, mẫu chia mẫu (dễ gặp phép chia có dư hoặc lẻ). Luôn thực hiện đúng quy tắc \textbf{Nhân đảo ngược}.
\end{itemize}
}
\end{tcolorbox}
\end{minipage}

\end{document}
